%\documentclass{Math_Note}
%\setlength{\parindent}{2em}

%\title{\textbf{Mathmatical Analysis Zorich Comprehensions Appendix}}
%\author{Buce-Ithon}
%\newdateformat{mydate}{\twodigit{\THEDAY}{ }\shortmonthname[\THEMONTH], \THEYEAR}
%\date{\today}

%\begin{document}
% Title page
%\maketitle

% Content
%\newpage
%\tableofcontents
%\newpage

% \setcounter{section}{-1}

\newpage
\section{Appendix}
\subsection{/Series/ A simple test for $n$th term of a series to approach $0$}
\textcolor{skycyan}{/Dev/ /Tec Tool/} \\
\fbox{\begin{minipage}{\textwidth}
Using Stirling’s formula, one may see at once that if $a_n = \frac{(2n)!}{4^n(n!)^2}$, then $a_n$ is of the order of $\frac{1}{\sqrt{n}}$, and one may conclude 
from the alternating series test that the series $\sum(-1)^n a_n$ is conditionally convergent. \\
At an elementary level, however, the convergence of the latter series may be a little more difficult to obtain. Since $\frac{a_{n+1}}{a_n} = \frac{(2n + 1)}{(2n + 2)} 
< 1$ for each $n$, it is clear that the sequence $\{a_n\}$ is decreasing, but it is not immediately obvious within the environment of a typical calculus course that 
$a_n \to 0$ as $n \to \infty$. For this purpose, one might use the following simple result which takes a leaf out of the theory of infinite products: 
\newline\newline
\fbox{\begin{minipage}{0.95\textwidth}
Lemma. Suppose $(a_n)$ is a decreasing sequence of positive numbers and for each natural number $n$, define $b_n = 1 - \frac{a_{n+1}}{a_n}$. Then the sequence 
$\{a_n\}$ converges to zero if and only if the series $\sum b_n$ diverges.
\end{minipage}}
\newline\newline
\fbox{\begin{minipage}{0.95\textwidth}
Proof. We note first that unless $b_n \to 0$ as $n \to \infty$, both of the series $\sum b_n$ and $\sum \log(1 - b_n)$ diverge. On the other hand, if 
$b_n \to 0$, then $b_n/(-\log(1 - b_n)) \to 1$ as $n \to \infty$, and it follows from the limit comparison test that $\sum b_n$ diverges if and only if 
$\sum \log(1 - b_n)$ diverges. We note also that since $0 < b_n < 1$, we have $\log(1 - b_n) < 0$ for each $n$.
\end{minipage}}
\newline\newline
Now since $1 - b_n = a_{n+1}/a_n$ for every $n$, it is clear that $a_n = a_1(1 - b_1)(1 - b_2)(1 - b_3) \cdots (1 - b_{n-1})$ for each $n \geq 2$, and we 
therefore conclude that $a_n \to 0$ iff $\log a_n \to -\infty$ iff $\log a_1 + \sum_{i=1}^n \log(1 - b_i) \to -\infty$ iff $\sum \log(1 - b_n)$ diverges iff 
$\sum b_n$ diverges. \\
Returning now to the above example, we see that $b_n = 1/(2n + 2)$ for each $n$, and the obvious divergence of $\sum b_n$ implies that $a_n \to 0$. The same 
technique gives an easy proof of the convergence of such series as $\sum((-1)^nn^n/e^{n}n!)$, and the series $\sum(\frac{\alpha}{n})$ of binomial coefficients 
with $\alpha > -1$.
\end{minipage}}
\textit{From: American Mathematical Monthly, Vol. 85, No. 10 (Page 942)}

%\end{document}